\section*{Introduction}
\addcontentsline{toc}{section}{Introduction}
\label{sec:introduction}

Fourteen-year-old Sewell Setzer III did not encounter a foundation model in a laboratory. He encountered a deployed relationship. Character.AI supplied a consumer application, persistent personas, memory, an interface, access rules, engagement design, warnings, and safety controls. Users could create characters and influence conversations. The complaint in \emph{Garcia v. Character Technologies, Inc.} also named the company's founders and Google, alleging different contributions of personnel, technology, knowledge, and resources. After months of interaction with a character modeled on a fictional television figure, Sewell died by suicide. His mother alleged wrongful death, negligence, product defects, failure to warn, deceptive practices, and related claims. At the pleading stage, the district court held that the application was plausibly alleged to be a product and declined, on that record, to classify the generated output as protected speech. The action later settled without a finding of defect, causation, or liability.\lrfootnote{See \casecite{garcia2025}; \casecite{garciaDismissal2026}; \casecite{garciaLien2026}. This Article states disputed facts as allegations and uses the published Rule 12 order for its classification analysis, not as a merits judgment.}

The case exposes a problem more durable than either procedural ruling. An injured person can observe the interface and the output. The conduct relevant to ordinary law lies elsewhere: which model and version ran; who chose the relational objective, system instructions, memory, age gate, crisis response, and engagement metric; what evaluations or complaints arrived; who could change or stop the service; and which enterprise retained each record. Those facts may be divided among a model provider, application company, infrastructure supplier, deploying institution, and individual decisionmakers. To plead reasonable conduct against the correct actor, the claimant needs the configuration and control record. To obtain that record, the claimant must first identify an actor who possesses it. The merits standard exists; access to the facts that feed it does not.

That circularity is the responsibility problem this Article addresses. Current debate often asks whether tort law can accommodate AI, whether software is a product, or whether an artificial agent could become a legal person. Each question matters in the case that actually presents it. None supplies an operating method for divided deployments. A reasonableness standard cannot evaluate conduct until the relevant conduct and actor are visible. A product classification does not reveal which company selected a dangerous harness. Artificial personhood would relocate the dispute to an entity with no independent assets, records, governance, or practical susceptibility to judgment.

Providers can intensify the problem through a rhetorical switch. Before injury, an application may be a companion that understands, remembers, plans, or acts. After injury, the same application may become a probabilistic model whose precise output no employee selected. The first description can cultivate use and reliance; the second can make responsibility appear to dissolve. Law should accept the factual content of both descriptions without accepting the transfer implicit in either. Anthropomorphic design is evidence of intended use, foreseeable exposure, and feasible precaution. Stochastic generation is evidence about the mechanism of control. Marketing cannot create an artificial legal subject, and token-level unpredictability cannot erase enterprise control over deployment.

The Article's central claim is that AI responsibility should be reconstructed around \emph{deployment control}. A model becomes a legally consequential system through choices about purpose, users, data, instructions, retrieval, memory, tools, credentials, monitoring, escalation, updates, and shutdown. Those choices form a control structure even where no person selects the final output. This Article represents that structure through a \emph{deployment-control cascade}. The cascade traces authority and resources downward, operational effects outward, and feedback upward. At each level it asks who set a constraint, what the actor believed about the system and environment, which signal could correct that belief, and who could pause, roll back, restrict, or terminate operation.

The cascade is adapted narrowly from systems-safety analysis, but its output is ordinary legal evidence. The actor who chose a control action may have authorized the relevant conduct. A risk report or repeated complaint may establish notice. A process model can show what the enterprise knew or should have known about users and safeguards. Tested age gating, a crisis interruption, or a rollback mechanism can bear on feasible alternative design. Version and incident records can prove or rebut a causal path. The method does not decide duty, breach, defect, causation, or damages. It makes the deployment facts available to the doctrine that does.

Relational systems demonstrate the payoff. Controlled and observational research now identifies provider-governed variables---relationship-seeking behavior, repeated exposure, continuity, memory, and intervention---that can change user outcomes. In a preregistered four-week study, relationship-seeking behavior increased separation distress by 6.04 percentage points and produced one additional dependency-risk profile for every eleven people exposed when combined with emotional conversation, compared with the most functional baseline; month-long exposure produced no discernible improvement in psychosocial health.\lrfootnote{See \bluecite{kirkSteering2026}. The study supports causal propositions about its controlled exposures, not population incidence or tort causation in \emph{Garcia}.} Other studies map experienced companion loss, emotional dependence, and recurring harmful interaction patterns.\lrfootnote{See \bluecite{defreitasReplika2025}; \artcite{banksDeletion2024}; \artcite{laestadiusTooHuman2024}; \bluecite{zhangDarkSide2025}.} These findings belong not only in a descriptive account of companion use. They identify variables for the cascade's process model, foreseeability inquiry, monitoring design, and alternatives analysis.

Making control visible permits a responsibility rule with two layers. The external layer gives the claimant a front door. For a designated high-impact deployment, every enterprise that deploys the system in its activity and possesses material control over the risk at issue should owe a nondelegable deployment duty and serve as a proper initial responsibility center. After a claimant pleads a recognized injury and a plausible material nexus to the deployment, that enterprise should identify the relevant entities and produce the proportionate, nonprivileged deployment and incident record. Multiple enterprises can occupy parallel lanes. Contribution, indemnity, individual fault, and ultimate allocation follow the evidence.

This is primary recourse without automatic recovery. The claimant retains every element imposed by the governing cause of action, including legally cognizable harm, breach or defect, causation, and damages. Jurisdiction, immunity, privilege, proportionality, and substantive defenses remain in place. The rule changes sequence: a company cannot require the claimant to reconstruct its control architecture before producing the record of that architecture, and it cannot defeat attribution merely because no employee composed the precise output.

The internal layer is a mandatory \emph{responsibility charter} adopted before release. The charter names risk owners, fixes deployment approval, records unresolved dissent, protects escalation, assigns tested stop and rollback authority, preserves configuration and incident evidence, and establishes indemnification and recourse. Its purpose is to align authority with accountability. The engineer who documented a danger and used the authorized channel should not occupy the same position as the manager who suppressed the report or the employee who knowingly removed a constraint. Externally, the enterprise remains answerable. Internally, the charter protects the reporter and preserves recourse against concealment, falsification, retaliation, and unauthorized bypass.

The architecture survives the strongest objections because its claim is deliberately exact. Kenneth Abraham and Catherine Sharkey show that objective reasonableness and conventional tort doctrine can reach many first-generation AI disputes without decoding every internal model state.\lrfootnote{See \artcite{abrahamSharkey2027}.} The proposal accepts their merits premise and addresses the identification-and-proof loop that precedes it. Mark Geistfeld shows how expanded liability and duplicative insurance can increase the price paid by the class liability is meant to protect.\lrfootnote{See \artcite{geistfeldInsurability2026}.} The front door does not enlarge the substantive class of compensable losses, while the charter is designed to supply deployment-level records that can inform risk classification. Zachary Henderson argues that AI activities are heterogeneous and that some AI-caused losses should have no tort remedy.\lrfootnote{See \artcite{hendersonRobotLiability2026}.} A legislature may draw the substantive boundary exactly there and still ensure that claims within it are screened by their merits rather than by the claimant's ability to discover a hidden org chart. The objections narrow the trigger and production obligation; they reinforce the need to distinguish procedure, governance, and substantive liability.

The Article also develops a second rule for the settings in which automation threatens to become an excuse for institutional withdrawal. The \emph{anti-substitution principle} provides that where law independently imposes a duty of care, process, accommodation, professional judgment, or public service, the availability of AI does not by itself discharge or narrow that duty. Automation may perform part of the obligation when the governing law permits and when the deployed process preserves the protection the duty exists to secure. The principle is neither a ban on automation nor a general right to human contact. It is a rule against treating technological availability as legal performance.

Neither contribution depends on resolving machine consciousness. Absent positive law creating and equipping a new juridical status, an AI system is not an independent bearer of the claims, duties, powers, assets, procedural capacity, and liabilities necessary to judgment. The same system may nonetheless be classified as a product, service, instrument, communicative medium, or professional component because each classification answers a different doctrinal question. Stable legal subjects and variable functional objects allow courts to reason about new systems without either a universal ontology or a machine liability sink.

The sequence of the Article reflects that priority. Part I introduces the deployment-control cascade and translates its variables into legal facts. Part II supplies the compact jurisprudential baseline that keeps responsibility attached to recognized persons while allowing functional classification to vary. Part III uses relational harm, copyright, and anthropomorphic marketing to show which deployment facts classification brings into view. Part IV constructs the enterprise-facing front door, tests it against the leading objections, applies the full rule to \emph{Garcia}, and develops the mandatory charter and production obligation. Part V gives the anti-substitution principle its own doctrinal form. Part VI addresses capability stratification and the boundary between lawful government procurement and coercive retaliation. The Conclusion leaves one organizing proposition: probabilistic output does not make control disappear; it changes where law must look for it. An unnumbered source note identifies the official public account and author translations used for the Hangzhou comparison.
